\subsection{Trace of a quadratic form with respect to another}

In this section, E will denote a real vector space and Q, H two positive quadratic forms on E. There exist two symmetric bilinear forms \((x, y) \mapsto \langle x | y \rangle_{\mathbf{Q}}\) and \((x, y) \mapsto \langle x | y \rangle_{\mathbf{H}}\) on \(E \times E\) , characterized by

\[
\mathrm{Q} (x) = \langle x | x \rangle_ {\mathrm{Q}}, \quad \mathrm{H} (x) = \langle x | x \rangle_ {\mathrm{H}}
\]

for all \(x \in \mathbf{E}\) .

We call the trace of Q with respect to H, and write \(\mathrm{Tr}(\mathrm{Q}/\mathrm{H})\) , a real positive number, finite or not, defined as follows :

\begin{enumerate}
\def\labelenumi{\alph{enumi})}
\tightlist
\item
  If there exists \(x \in \mathbf{E}\) with \(\mathbf{H}(x) = 0\) and \(\mathbf{Q}(x) \neq 0\) , we put \(\mathrm{Tr}(\mathbf{Q} / \mathbf{H}) = +\infty\) .
\item
  Otherwise, \(\mathrm{Tr}(\mathbf{Q} / \mathbf{H})\) is the upper bound of the set of all numbers of the form \(\sum_{i=1}^{m} \mathbf{Q}(x_i)\) where \((x_1, \ldots, x_m)\) range over the set of finite sequences of elements of \(\mathbf{E}\) such that \(\langle x_i | x_j \rangle_{\mathbf{H}} = \delta_{ij}\) (Kronecker's symbol).
\end{enumerate}

Remarks. --- 1) For every subspace F of E, let \(\mathrm{Q}_{\mathrm{F}}\) denote the restriction of Q to F and \(\mathrm{H}_{\mathrm{F}}\) that of H. We have \(\mathrm{Tr}(\mathrm{Q}_{\mathrm{F}} / \mathrm{H}_{\mathrm{F}}) \leqslant \mathrm{Tr}(\mathrm{Q} / \mathrm{H})\) and \(\mathrm{Tr}(\mathrm{Q} / \mathrm{H})\) is the upper bound of the set of all numbers \(\mathrm{Tr}(\mathrm{Q}_{\mathrm{F}} / \mathrm{H}_{\mathrm{F}})\) where F ranges over the family of all finite dimensional vector subspaces of E.

\begin{enumerate}
\def\labelenumi{\arabic{enumi})}
\setcounter{enumi}{1}
\tightlist
\item
  Let \(\mathbf{E}_1\) be a real vector space, \(\mathbf{Q}_1\) and \(\mathbf{H}_1\) two positive quadratic forms on \(\mathbf{E}_1\) and \(\pi: \mathbf{E} \to \mathbf{E}_1\) a surjective linear mapping. If \(\mathbf{Q} = \mathbf{Q}_1 \circ \pi\) and \(\mathbf{H} = \mathbf{H}_1 \circ \pi\) , then \(\operatorname{Tr}(\mathbf{Q}/\mathbf{H}) = \operatorname{Tr}(\mathbf{Q}_1/\mathbf{H}_1)\) .
\end{enumerate}

PROPOSITION 15. --- Suppose that there exists a real hilbertian space structure on E such that \(\mathbf{H}(x) = \|x\|^2\) for all \(x \in E\) . For \(\operatorname{Tr}(\mathbf{Q}/\mathbf{H})\) to be finite, it is necessary and sufficient that there exists a continuous and positive endomorphism u of E with finite trace, such that \(\mathbf{Q}(x) = \langle x|u(x) \rangle\) for all \(x \in E\) ; this endomorphism u is unique, and we have

\[
\operatorname{Tr} (u) = \operatorname{Tr} (\mathrm{Q} / \mathrm{H}) = \sum_ {i \in \mathrm{I}} \mathrm{Q} \left(e _ {i}\right)\tag{44}
\]

for every orthonormal bases \((e_i)_{i\in \mathbf{I}}\) of E.

Suppose that \(\operatorname{Tr}(\mathbf{Q} / \mathbf{H})\) is finite. For every \(x\in E\) of norm 1, we have \(\mathrm{H}(x) = 1\) , hence \(\mathrm{Q}(x)\leqslant \mathrm{Tr}(\mathrm{Q} / \mathrm{H})\) . Therefore, \(\mathrm{Q}(x)\leqslant \mathrm{Tr}(\mathrm{Q} / \mathrm{H}).\| x\|^2\) for all \(x\in E\) , and

\[
\left| \langle x | y \rangle_ {\mathrm{Q}} \right| \leqslant \mathrm{Q} (x) ^ {1 / 2} \mathrm{Q} (y) ^ {1 / 2} \leqslant \operatorname{Tr} (\mathrm{Q} / \mathrm{H}). \| x \|. \| y \|
\]

by the Cauchy-Schwarz inequality. Consequently, the bilinear form \((x,y)\mapsto\langle x|y\rangle_{\mathbf{Q}}\) on \(E\times E\) is continuous. There exists (V, p.~16, cor. 2) a mapping \(u\in\mathcal{L}(E)\) such that \(\langle x|y\rangle_{\mathbf{Q}}=\langle x|u(y)\rangle\) . We have \(\langle x|y\rangle_{\mathbf{Q}}=\langle y|x\rangle_{\mathbf{Q}}\) for x, y in E, hence u is hermitian; and \(\langle x|u(x)\rangle=\mathbf{Q}(x)\geqslant0\) , hence u is positive.

Conversely, let \(u\) be a continuous and positive endomorphism of \(E\) such that \(Q(x) = \langle x|u(x) \rangle\) for all \(x \in E\) . Then

\[
\langle x | u (y) \rangle = \frac {1}{2} (\mathrm{Q} (x + y) - \mathrm{Q} (x) - \mathrm{Q} (y)) = \langle x | y \rangle_ {\mathrm{Q}},
\]

which gives the uniqueness of u. By formula (24') (V, p.~50), we get

\[
\operatorname{Tr} (u) = \sup _ {x _ {1}, \dots , x _ {m}} \sum_ {i = 1} ^ {m} \left\langle x _ {i} \mid u \left(x _ {i}\right) \right\rangle = \sup _ {x _ {1}, \dots , x _ {m}} \sum_ {i = 1} ^ {m} Q \left(x _ {i}\right),
\]

where \((x_{1},\ldots ,x_{m})\) range over the set of all finite orthonormal sequences of elements of E. By the definition of \(\mathrm{Tr}(\mathrm{Q / H})\) , we get \(\mathrm{Tr}(u) = \mathrm{Tr}(\mathrm{Q / H})\) . Finally, for every orthonormal basis \((e_i)_{i\in \mathbf{I}}\) of E, we have \(\mathrm{Tr}(u) = \sum_{i\in \mathbf{I}}\langle e_i|u(e_i)\rangle\) by formula (25) of V, p.~50, hence \(\mathrm{Tr}(u) = \sum_{i\in \mathbf{I}}\mathrm{Q}(e_i)\) . Q.E.D.

Remarks. --- 3) Let E and F be two hilbertian spaces and v a linear, not necessarily continuous mapping from E into F. Let \(\mathbf{H}(x)=\|x\|^{2}\) and \(\mathbf{Q}(x)=\|v(x)\|^{2}\) for all \(x\in E\) . It follows from prop. 15 that v is a Hilbert-Schmidt mapping if and only if \(\operatorname{Tr}(\mathbf{Q}/\mathbf{H})\) is finite, and then \(\operatorname{Tr}(\mathbf{Q}/\mathbf{H})=\|v\|_{2}^{2}\) .

\begin{enumerate}
\def\labelenumi{\arabic{enumi})}
\setcounter{enumi}{3}
\tightlist
\item
  Suppose E is finite dimensional. When the quadratic form H is invertible, prop. 15 applies. Let \((e_{1}, \ldots, e_{n})\) be a basis of E. Put \(q_{ij} = \langle e_{i}|e_{j} \rangle_{Q}\) and \(h_{ij} = \langle e_{i}|e_{j} \rangle_{H}\) and introduce the matrices \(q = (q_{ij})\) and \(h = (h_{ij})\) . Let u be an endomorphism of E such that \(Q(x) = \langle x|u(x) \rangle_{H}\) for all \(x \in E\) . We have
\end{enumerate}

\[
\langle x | y \rangle_ {\mathrm{Q}} = \langle x | u (y) \rangle_ {\mathrm{H}} \quad (x, y \in \mathrm{E}),
\]

and hence the matrix of \(u\) with respect to the basis \((e_1, \ldots, e_n)\) of \(E\) is equal to \(h^{-1}q\) . By prop. 15, we have

\[
\operatorname{Tr} (\mathrm{Q} / \mathrm{H}) = \operatorname{Tr} (h ^ {- 1} q) = \operatorname{Tr} (q h ^ {- 1}).\tag{45}
\]

If the basis \((e_1, \dots, e_n)\) is orthonormal for \(\mathbf{H}\) , then \(h\) is the unit matrix of order \(n\) , and we get

\[
\operatorname{Tr} (\mathrm{Q} / \mathrm{H}) = \operatorname{Tr} (q) = \sum_ {i = 1} ^ {n} \mathrm{Q} (e _ {i});
\]

so that we get formula (44) in this case.

Now suppose that the quadratic form \(\mathbf{H}\) is not invertible. Let \(\mathbf{N}\) be the kernel of \(\mathbf{H}\) , and let \(\pi\) be the canonical mapping from \(\mathrm{E}\) onto \(\mathrm{E} / \mathrm{N}\) . There exists an invertible quadratic form \(\mathrm{H}_1\) on \(\mathrm{E} / \mathrm{N}\) such that \(\mathrm{H} = \mathrm{H}_1 \circ \pi\) . Let \((e_1, \dots, e_n)\) be a sequence of elements of \(\mathrm{E}\) such that the sequence \((\pi(e_1), \dots, \pi(e_m))\) is a basis of \(\mathrm{E} / \mathrm{N}\) , which is orthonormal for \(\mathrm{H}_1\) . Let \((e_1, \dots, e_n)\) be a basis of \(\mathbf{N}\) . Then \((e_1, \dots, e_n)\) is a basis of \(\mathrm{E}\) and we have

\[
\mathrm{H} \left(\xi_ {1} e _ {1} + \dots + \xi_ {n} e _ {n}\right) = \xi_ {1} ^ {2} + \dots + \xi_ {m} ^ {2}\tag{46}
\]

for all real numbers \(\xi_1, \ldots, \xi_n\) .

Suppose that for all \(x \in \mathbf{E}\) , the relation \(\mathrm{H}(x) = 0\) implies \(\mathrm{Q}(x) = 0\) ; in other words, suppose that there exists a quadratic form \(\mathbf{Q}_1\) on \(\mathrm{E} / \mathrm{N}\) such that \(\mathrm{Q} = \mathrm{Q}_1 \circ \pi\) . By remark 2 and prop. 15, we have,

\[
\operatorname{Tr} (\mathrm{Q} / \mathrm{H}) = \mathrm{Q} (e _ {1}) + \dots + \mathrm{Q} (e _ {m}).\tag{47}
\]

